\documentclass[10pt,a4paper]{ltjsarticle}
\usepackage[margin=14mm,top=10mm,bottom=10mm]{geometry}
\usepackage{amsmath,amssymb}
\usepackage{enumitem}
\usepackage[hidelinks]{hyperref}
\setlist{nosep,leftmargin=1.6em}
\pagestyle{empty}
\setlength{\parskip}{0.2em}
\newcommand{\sect}[1]{\par\vspace{0.45em}\noindent{\bfseries\sffamily #1}\par\vspace{0.1em}}
\newcommand{\Cmcf}[2]{\mathcal{C}^{\mathrm{mcf}}_{#1,#2}}

\begin{document}

\begin{center}
{\large\bfseries 固定有限半群型付けの下での多重文脈自由言語の分布学習：\\局所群と特徴データの長さ}\\[0.3em]
{\small Distributional Learning of Multiple Context-Free Languages under Fixed Finite-Semigroup Typing:\\
Local Groups and the Length of Characteristic Data}\\[0.3em]
栗山 貴之（独立研究者）\quad 概要（2026年11月，v2）
\end{center}

\sect{1. 設定}
有限半群の上への型付け $s:\Sigma^+\twoheadrightarrow\mathsf S$ を，データを見る前に固定します．
タプル成分は $s$-値が一致するときにのみ比較し，標的クラス $\Cmcf{f}{s}$ は，この制限の下で
タプル代入可能なファンアウト $f$ の MCFL からなります．再構成は，吉仲先生の
Const/Link/Comp の作用素に，Link の条件として型の一致を加えたものです．
質的な同定は先生の枠組みから従います．本稿の主題は，特徴標本の\textbf{符号化総長}を，
標的の既約な良い表示 $G$ の大きさ $\|G\|$ と規則の厚さ $\bar\theta_G$ で測ることです．

\sect{2. 主結果}
\begin{itemize}
\item \textbf{定理 A（局所群）．}$s$ の正の像が局所群（すべてのべき等元 $e$ で $e\mathsf Se$ が群）ならば，
固定したすべてのファンアウトとランクの上限で，すべての $L\in\Cmcf{f}{s}$ は，すべての既約な良い表示 $G$ について
長さ $O(\|G\|^2\bar\theta_G)$ の特徴データをもちます．
\item \textbf{定理 B（線形因子）．}$4f+1$ 文字以上（$f=1$ では2文字）のアルファベット上で，すべての型付けについて，
最小に近い表示をもつ有限の標的で，特徴標本がすべて長さ $\Omega(\|G\|\bar\theta_G)$ をもつものがあります．
自明な型付けでは，先生の上界 $O(\|G\|\bar\theta_G)$ が定数倍を除いて達成されます．
\item \textbf{定理 C（ファンアウト1での境界）．}$f=1$，ランク2以上で，有限半群の擬多様体が一様に多項式の
特徴データ長をもつのは，それが $\mathbf{LG}$（局所群）に含まれるとき，かつそのときに限ります．
外側では，二元半束 $U_1$ の上への型付けについて，大きさと厚さが $O(n)$ の表示をもつ標的が長さ $3\cdot2^n-1$ 以上の
データを要します（クラス全体で特徴標本をもつすべての集合駆動型作用素について）．
ファンアウト1では $\mathbf{LG}$ がちょうど IPTtD 可能な境界です．
\end{itemize}
型なしのクラスとの比較として，$\mathbb Z_2$-型付けの下でクラスに属し，型なしの作用素では過汎化する
非文脈自由な 2-MCFL の例も与えています．

\sect{3. 証明の要点}
\begin{itemize}
\item ファンアウト $f\ge2$ では，部分木が最大 $f$ 個のブロックを異なる根成分に与え，他の部分木のブロックと交互に並びます．
非可換群では置換の効果がこの順序に依存するので，すべての根成分の順序付き群積を同時に保つ置換を
\textbf{所有者ごとに同時に選ぶ}補題（Ramsey 型の順序の均質化と，接頭辞積の鳩の巣論法）を用います．
\item 局所群は Rees の座標と滑動因子の群寄与で表し，同じ選択に帰着させます．下界は稀な型（その型の唯一の語）を強制する議論によります．
\end{itemize}

\sect{4. ファンアウト2以上（未解決）}
定理 C(ii) の倍増族は，二つの穴をもつ冪の移送により，任意の固定した型付けの下でファンアウト2に持ち越されません．
\textbf{$\mathbf{LG}$ がファンアウト $f\ge2$ でも擬多様体の境界かどうか}は未解決問題として残しています．
（原稿外の探索では，二文字の内容型付けはファンアウト2で多項式になることまで分かっています．）

\sect{5. 先生にうかがいたいこと}
\begin{enumerate}
\item \textbf{型付きの再構成と，先生の 2015 年の枠組みとの境目．}本稿の再構成は，先生の有限シグネチャの枠組み
（Yoshinaka 2015）の MCFG シグネチャで，タプルのソートを $s$-型で細分したものとして説明しています（同論文 Section~8 を参照）．
先生から見て，型付きのクラスとその質的な同定は「既にその枠組みに含まれている」ものでしょうか．
それとも型付けを固定したバイアスとして導入すること自体に意味があるとお考えでしょうか．
\item \textbf{特徴標本の下界は既知でしょうか．}定理 B は，自明な型付けで先生の上界 $|P|\theta_G$（Yoshinaka 2011）が
定数倍を除いて達成されること，つまり大きさと厚さの積が避けられないことを示しています．
この下界，あるいは厚さが本質的に必要であることは，既に知られていたでしょうか．
\item \textbf{ファンアウト2以上の境界についてのご直観．}先生の 2011 年の論文には，$1$-次元では代入可能でも $2$-次元では代入可能でない例があります．
ファンアウト2の代入可能性は，型付けで強制される長い語を壊す方向に働く（境界が動く）と見えますか．
それとも $U_1$ 型の障害がどこかで再び現れる（境界は $\mathbf{LG}$ のまま）と見えますか．
\end{enumerate}

\vspace{0.3em}\noindent{\footnotesize 英語原稿（44ページ）：\url{https://github.com/growupkuriyama-hub/Papers/blob/main/02-v2-lg-fanout-boundary/main.pdf}\\
日本語版（49ページ）：\url{https://github.com/growupkuriyama-hub/Papers/blob/main/02-v2-lg-fanout-boundary/japanese/main_JP.pdf}}

\end{document}
